% Prior work on proof simplification~\citep{ahuja2024improver} focuses on agentic scaffolding around API-only LLMs such as GPT-4o. While these methods can shorten human-written Lean proofs, they are ineffective at simplifying the long proofs generated by SoTA RL-trained LLM provers such as Seed-Prover and Goedel-Prover-V2~\citep{lin2025goedel2}, precisely the setting where simplification is most valuable. A natural alternative is to finetune LLMs directly for proof simplification, but progress in this direction is limited by the lack of suitable training data, namely aligned pairs of proofs before and after simplification.
%
% The notion of \textit{simple proof} can be specified in different ways~\citep{kinyon-proof-simp}. In this work, we choose a natural definition of simplicity: the length of a proof, calculated as the number of Lean tokens in it. Finding the shortest proof is in general a challenging problem: intuitively, it is exponentially more difficult than finding \textit{any} valid proof, and no efficient symbolic techniques exist for shortening proofs in Lean.

By focusing on proof simplicity as the target metric, we developed \emph{ProofOptimizer}, an
LLM-based system for simplifying long and convoluted proofs in Lean.
ProofOptimizer integrates three components: (i) a symbolic Lean linter that identifies and
removes redundant steps, (ii) a 7B parameter language model finetuned specifically
for proof simplification, and (iii) an iterative inference-time algorithm for progressively
shortening proofs. Given an input proof, the Lean linter first eliminates the most
obvious redundancies. The language model then generates multiple candidate simplifications, and the iterative algorithm repeatedly applies the model to the currently shortest proof, further reducing its length. Training follows two paradigms. In expert iteration, the model proposes simplifications that are verified by Lean and incorporated into the training data for supervised finetuning. In reinforcement learning, proof length and correctness serve as the reward signal. Both approaches enable continual improvement without requiring any human-annotated simplification data.

% Therefore, we aim to mitigate the problem of suboptimal long proofs and we develop ProofOptimizer: a LLM-based approach for optimizing the length of Lean proofs. ProofOptimizer combines a symbolic Lean linter, a 7B parameter language model fine-tuned by us specifically for the proof simplification task, and an iterative proof shortening inference-time algorithm.

We evaluated ProofOptimizer on long proofs generated by state-of-the-art neural theorem provers.
Our final models achieved decent results (Fig. \ref{fig:example}), shortening MiniF2F proofs by an average of 63\% in a single shot
and PutnamBench proofs by 26\% with 32 attempts, substantially outperforming Gemini-2.5-Pro (Sec.~\ref{sec:expit-rl}).
At test-time, iterative shortening helped us achieve per-proof average reductions of 87\% (MiniF2F) and 57\% (PutnamBench)
and reduced the length of three out of four Seed-Prover IMO 2025 proofs by more than half.

Following the second lesson of this thesis, we did further analysis and found two downstream benefits of proof shortening.
Training our base model on shortened proofs led to 2\% better performance on miniF2F relative to training on unshortened proofs (Sec.~\ref{subsec:training-simplified}).
Also, shortening proofs often decreased their execution time, with 28\% of proofs showing at least a 1.5x speedup after shortening (Sec.~\ref{subsec:proof-execution-time}).

\begingroup
\setlength{\abovecaptionskip}{6pt}
\setlength{\belowcaptionskip}{0pt}
\begin{figure}[H]
    \centering
    \includegraphics[width=0.95\linewidth]{figures/main_fig_large.pdf}
    \caption{Overview of our pipeline. ProofOptimizer reduces the shortest generated proof of a Putnam problem from 1097 to 76 tokens.}
    \label{fig:example}
\end{figure}
\endgroup
